\scanpage{9}
\mspage{5}
% archive sheet number: 8

$\unsure{\Phi|\mathscr F},\Phi'$, et une seule). \deleted{Ainsi.} Quand on se limite aux polyfigures \inserted{(généralisées)}, \unsure{irrédundantes}, \deleted{\unsure{double|donc}} l'introduction des $I_{<a}$, $I_{>a}$ on aura $\delta\Phi=\delta\Phi'$, alors que ce ne sera pas le cas pour le formalisme avec les polyfigures non généralisées, où les extrémités \inserted{de $L$} viennent nous emmerder.

Introduction de nouveaux éléments, pour former $\widehat{\unsure{\mathscr M|\mathcal M}}=\widehat{\mathscr M}_0\amalg\widehat{\mathscr M}_1$.
\[
\begin{aligned}
\widehat{\mathscr M}_0&=\mathscr M_0=L\\
\widehat{\mathscr M}_1&=\mathscr M_1\amalg\mathscr M_1^-\amalg\mathscr M_1^+\\
\mathscr M_1^-&=\{I_{<a}\mid a\in L,\ a\text{ non minimal dans }L\}\\
\mathscr M_1^+&=\{I_{>a}\mid a\in L,\ a\text{ non maximal dans }L\}
\end{aligned}
\]
$I_{<a}$, $I_{>a}$ \deleted{\unsure{pour}} sont ici des symboles graphiques, auxquels pour l'instant on n'essaie de donner une interprétation. Heuristiquement, \deleted{\unsure{elle fig}} $I_{<a}$ figure le demi-intervalle formé des $x\in L$, $x<a$, ou mieux, les $X\in\mathscr M$ tels que $X<a$, et symétrique pour $\unsure{L_{>a}|I_{>a}}$.

On pose
\[
\begin{aligned}
\partial I_{\unsure{<a|\leq a}}&=\delta I_{<a}=\deleted{\unsure{\partial I_{>b}=\delta I_{<a}}}\unsure{=}\{a\}\\
|I_{<a}|&\overset{\mathrm{d\acute ef}}{=}L_{\leq a}=\{x\in L\mid x\leq a\}\\
|I_{<a}|^\circ&\overset{\mathrm{d\acute ef}}{=}\deleted{\unsure{L_{<a}}}|I_{<a}|\setminus\partial I_{<a}=L_{<a}=\{x\in L\mid x<a\}
\end{aligned}
\]
et dualement pour $I_{>a}$. Cela étend $\unsure{\unlhd|\trianglelefteq}$ et $\triangleleft$ à $\widehat{\mathscr M}$. On \deleted{\unsure{intro}} prolonge de $=$ $\overset{\circ}{\ll}$, en écrivant les nouveaux cas de $\overset{\circ}{\ll}$ impliquant $I_{<a}$
